% Generated by make_size_class_copilot_pdf.py from copilot-prompts.md P11; edit those, not this file.
\documentclass[11pt]{article}
\usepackage[letterpaper, top=0.5in, bottom=0.5in, left=0.5in, right=0.5in, headheight=14pt, includeheadfoot]{geometry}
\usepackage[T1]{fontenc}
\usepackage{lmodern}
\usepackage{microtype}
\usepackage{enumitem}
\usepackage{titlesec}
\usepackage{multicol}
\usepackage{fvextra}
\usepackage{fancyhdr}
\usepackage{lastpage}
\usepackage[hidelinks, pdftitle={Open Orders Report: Build the Size Class File in Copilot (P11)}, pdfauthor={Gust F.}]{hyperref}

\pagestyle{fancy}
\fancyhf{}
\fancyhead[L]{\small Strategy}
\fancyhead[R]{\small Rev.~2 \quad 2026-10-15}
\fancyfoot[L]{\small Gust F.}
\fancyfoot[C]{\small Page \thepage\ of \pageref{LastPage}}
\renewcommand{\headrulewidth}{0.4pt}

\titleformat{\section}{\small\bfseries}{Step \thesection.}{0.5em}{}
\titlespacing*{\section}{0pt}{0.35em}{0.1em}
\setlist[enumerate]{leftmargin=1.8em, itemsep=0em, topsep=0.1em, label=(\alph*), ref=\alph*}
\setlength{\parindent}{0pt}
\setlength{\parskip}{0pt}
\setlength{\columnsep}{0.25in}

\newcommand{\menu}[1]{\textit{#1}}
\newcommand{\gt}{~>\ }
% Prompt text, copied word for word from copilot-prompts.md P11 (tabs shown as <TAB>): framed, wrapped, ready to copy and paste.
\DefineVerbatimEnvironment{prompt}{Verbatim}{breaklines, breakanywhere=false, frame=single, framesep=0.3em,
  fontsize=\fontsize{6.1}{6.8}\selectfont, fontfamily=lmss, breaksymbolleft={}, breakindent=0.8em}
\newcommand{\checknote}[1]{\par\vspace{0.15em}{\footnotesize\textbf{Check.} #1}\par}

\begin{document}

\begin{center}
  {\large\bfseries Build the Size Class File in Copilot (P11)}\\[0.25em]
  {\small From the PowerPoint slide to the workbook Config \texttt{SizeClassFile} points at}
\end{center}

\vspace{-0.7em}
\noindent\rule{\linewidth}{0.8pt}

{\footnotesize\textbf{Instructions.} Do this first: it is the most important Copilot prompt. The macro builds nothing without a size class file, and a model missing from it lands in UNMAPPED. The slide has the latest classes, but not in the macro's layout or MODEL\_GROUP spelling. Run Steps~\ref{sec:s1} to \ref{sec:s3} \textbf{in one Copilot chat}, in order, attaching the file each one names; Step~\ref{sec:s4} is optional. Paste each boxed prompt as it is; where it shows <TAB>, a tab is meant. If Copilot loses the earlier table, start a new chat and paste it under the prompt.\par}

\vspace{-0.1em}
\noindent\rule{\linewidth}{0.4pt}

\begin{multicols}{2}\footnotesize
\section{Read the slide (attach the PowerPoint file)}\label{sec:s1}
\begin{prompt}
Attached is a PowerPoint file. One slide in it lists Bobcat machine models and the size class each belongs to. The layout may be a table, groups of models under class headings, or columns per class.

Task: build one table with one row per machine model on that slide.

How to fill each row:
- Name of machine: the model code only, as written on the slide, without the word Bobcat or the machine type. Examples: T76, S650, E35, T76-2, E50Z. If the slide writes a series as a suffix (for example "T76 R2" or "T76 R-Series"), keep it exactly as written; Step 2 fixes the spelling.
- Slide class: the class label exactly as the slide shows it for that model.
- Size Class: the slide class translated to exactly one of these names, nothing else:
  Track loaders (T...), by rated operating capacity in lb, as the business defines them:
    0<2000, 2000<2300, 2300<2600, 2600<3000, 3000<3500, 3500&Over
  Mini excavators (E...), by nominal tonnes, as the business defines them:
    >00-01, >01-02, >02-03, >03-04, >04-05, >05-06, >06-10
  Skid-steer loaders (S...), by frame size, kept exactly as the slide names them:
    EX Small, Small Frame, Med Frame, Large Frame
- Rated Operating Capacity: the capacity in lb if the slide shows it (for an excavator, its operating weight in lb if shown); otherwise leave it blank.
- HP: the horsepower if the slide shows it; otherwise leave it blank.
- Notes: write "CHECK: slide class does not match" if the slide label does not translate to exactly one of the names above (for example its range crosses a boundary), "CHECK: on a class boundary" if a capacity shown is within 100 lb of a track loader class edge, and "CHECK: listed twice" if a model appears more than once on the slide. Otherwise leave it blank.

Rules:
- Use only what is on the slide. Do not add models from memory. Do not look anything up.
- Do not explain. Reply only in the format below.
- If you cannot find such a slide, reply with one line: CANNOT: <reason>

Reply in exactly this format:
SLIDE: <slide number and title>
CLASS LABELS
<slide class label> -> <Size Class name, or NONE>
TABLE
Name of machine | Slide class | Size Class | Rated Operating Capacity | HP | Notes
<one line per model>
MODELS ON SLIDE: <number>
\end{prompt}
\checknote{Count the models on the slide and compare with MODELS ON SLIDE. Each CLASS LABELS line must point to the class you expect; if not, correct Copilot in the chat (e.g.\ \texttt{"Under 2,000" is 0<2000}) and ask it to redo the table.}

\section{Check against Order Intake (attach it)}\label{sec:s2}
\begin{prompt}
Attached is the Order Intake export. Use its first sheet. Keep the table from the slide.

Task: compare the slide table with the models in this export.
1. Take the export rows whose PRODUCT_TYPE_ABBR is CTL, MX or SSL (ignore capitals and extra spaces). List every different MODEL_GROUP among them.
2. For each MODEL_GROUP, find its row in the slide table: the same name ignoring capitals and spaces, or the same model written another way (for example "T76 R2" on the slide and "T76-2" in the export). When it is the same model written another way, change Name of machine in the slide table to the export's spelling and add "renamed from <slide name>" to Notes.
3. A MODEL_GROUP with no row in the slide table: add a row with Name of machine = the MODEL_GROUP exactly as in the export.
   - E followed by digits (a mini excavator): leave Size Class blank; Notes "CHECK: not on slide, class unknown". There is no rule for it: the class comes from the slide or the spec sheet.
   - Any other model: give your best Size Class only if you are sure of its rated operating capacity; Notes "CHECK: not on slide". If you are not sure, leave Size Class blank and write "CHECK: not on slide, class unknown".
4. Keep every slide row, including models the export does not have.

Rules:
- Read every row of the export. Do not explain. Reply only in the format below.
- If a column is missing or you cannot read the file, reply with one line: CANNOT: <reason>

Reply in exactly this format:
RENAMED
<slide name> -> <MODEL_GROUP>
ADDED (not on slide)
<PRODUCT_TYPE_ABBR> | <MODEL_GROUP> | <rows in export> | <Size Class or blank>
ON SLIDE, NOT IN EXPORT
<Name of machine>
TOTAL MODEL_GROUPS IN EXPORT: <number>
\end{prompt}
\checknote{Every blank Size Class needs one before Step~\ref{sec:s3}: find the capacity on the spec sheet and tell Copilot, e.g.\ \texttt{T870 is 3500\&Over}. A model left blank stays UNMAPPED.}

\section{The finished table (same chat)}\label{sec:s3}
\begin{prompt}
Now give the final size class table: the slide table after Step 2, with the corrections I gave you in this chat.

Rules:
- One row per Name of machine. No blank Size Class. Size Class must be exactly one of the 17 names from Step 1.
- Sort: track loaders (T...), then skid-steer loaders (S...), then mini excavators (E...), then anything else; inside each group by Size Class in the order listed in Step 1, then by Name of machine.
- Reply only with the table as tab-separated text, no other words, starting with this header row:
Name of machine <TAB> Rated Operating Capacity <TAB> HP <TAB> Size Class <TAB> Notes
- After the table, one line: ROWS: <number of model rows>
\end{prompt}
\vspace{0.1em}
{\footnotesize If your Copilot can create files, ask instead: \textit{Create an Excel workbook with this table on the first sheet, named Size Class, and Size Class formatted as text}; then go to step~(\ref{step:save}).}
\begin{enumerate}\footnotesize
  \item Open a new workbook; set column D (Size Class) to \menu{Home\gt Number Format\gt Text} before pasting.
  \item Paste into cell A1 of the \textbf{first sheet} (the macro reads only that one).
  \item Check that ROWS matches the rows pasted and no Size Class cell is empty.
  \item Save it as a local \texttt{.xlsx} (e.g.\ in \texttt{C:\textbackslash Temp\textbackslash oo-test}).\label{step:save}
  \item On Config, click \menu{Choose size class file\ldots}, pick it, then click \menu{Check setup}; fix anything it lists in the same chat and paste again.
\end{enumerate}

\section{Optional: what changed}\label{sec:s4}
{\footnotesize Attach \texttt{size-class-master.xlsx} to show what the slide changed.\par}
\begin{prompt}
Attached is the previous size class list (first sheet; columns Name of machine and Size Class). Compare it with the final table from Step 3.

Rules: compare names ignoring capitals and spaces. Do not explain. Reply only in the format below.

Reply in exactly this format:
CLASS CHANGED
<Name of machine> | <previous Size Class> -> <new Size Class>
NEW (not in the previous list)
<Name of machine> | <Size Class>
DROPPED (in the previous list only)
<Name of machine> | <Size Class>
UNCHANGED: <number>
\end{prompt}
\end{multicols}

\end{document}
